% ======================================================================
% 六、问题三：日内滚动购电与调整策略的建模与求解（7 面版）
% 对应仓库：Sen-illion/math-modeling @ task/integrate-q123
% 冻结口径：LA = 48h 展望 + 三段分源残差分位修正（q_L=0.60 一月定参）
%           正式窗 2025-02-01--12-31，2/1 0:00 SOC=6000 kWh
% 需用宏包：amsmath, amssymb, booktabs, graphicx, caption, enumitem
% 图文件：fig_q3_horizon.png, fig_q3_dispatch_0621.png
% ======================================================================
\section{问题三：日内滚动购电与调整策略的建模与求解}

\subsection{问题分析与建模思路}

问题三在问题二“日前一次决策”的基础上，把一天拆成四个决策时点：$0{:}00$ 制定全天计划购电 $G^{\mathrm{plan}}$，随后在 $6{:}00$、$12{:}00$、$18{:}00$ 获得更新的光伏预报后，对尚未执行的合同做增购或减购调整，最终执行 $G^{\mathrm{adj}}$。调整并非免费：比计划少买按 $0.5\pi$ 结算、多买按 $1.5\pi$ 结算，仍出现缺口则按 $5\pi$ 紧急购电。这一问的难点不在“再解一次优化”，而在三件事环环相扣，我们先把它们拆开。

\textbf{（1）信息层——预报如何进入决策。} 附件 3 只给出四个发布时点之后 $T{+}1,\dots,T{+}24$ 的整点光伏预报，既不是 $10$ 分钟粒度，也没有“发布当刻”的起点值。若按整点硬拼，夏季 $6{:}00/18{:}00$ 起点会被系统压成零、低估全天光伏；若误用未发布的次日预报又造成信息泄漏。因此第一件事是\textbf{无泄漏地把整点预报补成 $10$ 分钟序列}，否则后面所有优化都建立在错误输入上。

\textbf{（2）结算层——不对称惩罚决定“分位怎么取”。} 少买 $0.5\pi$、多买 $1.5\pi$、缺口 $5\pi$ 三种价格不对称，直接改变“买多少保险”的答案：两次调整之间\textbf{锁死}的时段下一时刻前无法再改、少买代价高，应按高覆盖率买足；下一时刻\textbf{还能改}的时段，偏差可被下一轮修正，取中位即可。这正是后文“三段分位”的由来。

\textbf{（3）物理层——真正的失效在晚间电量预算。} 储能可用区间只有 $[1200,10800]$ kWh、接口 $5000$ kW。诊断（见 6.8 节）显示，旧方案紧急购电 $88.8\%$ 来自\textbf{晚间储能电量耗尽}而非功率触顶，$20{:}00$ 单小时占 $47.4\%$，日末有 $106$ 天贴近下限。模型必须在优化阶段就为晚间留够电量，而不是事后用 $5\pi$ 补缺口。

综合三层约束，本问被定位为一个\textbf{带追索权（recourse）的多阶段随机线性规划}：第一阶段在预报不确定时先承诺合同，后续阶段看到新预报后再做有代价的修正。目标与约束关于购电量、充放电量都是线性的，不确定性只通过“预报残差”进入，故用“分位点预报”把随机性近似掉，得到可逐时段求解的\textbf{滚动时域线性规划}；不必上随机规划或元启发式，因为随机性来源单一、惩罚结构明确，分位修正已能把风险压到计划阶段。

\subsection{模型假设}

本问沿用前两问公共假设：充放电单程效率各为 $\eta=0.9$（交流侧口径）、储能运行窗 $[1200,10800]$ kWh、接口最大功率 $5000$ kW、不向外网倒送电、光伏过剩允许弃光；一天仍为 $144$ 个 $10$ 分钟时段，$\Delta t=1/6$ h。相对问题二另有三条约定：光伏决策只用附件 3 预报并线性插值到 $10$ 分钟；负荷 $0{:}00$ 沿用因果 Expanding XGBoost、$6/12/18$ 时点用已发生均值比例校正；不强制 $E_0=E_{24}$，储电量跨日传递，正式窗 $2$ 月 $1$ 日 $0{:}00$ 置 $6000$ kWh、$1$ 月仅用于预报与定参。新增符号 $G^{\mathrm{plan}},G^{\mathrm{adj}},\Delta^{\pm},G^{\mathrm{em}},q_L,q_O,q_E$ 在正文就近定义。

\subsection{结算规则与目标函数}

设一天 $144$ 段，$t=1,\dots,144$，交易电价 $\pi_t$。定义上调量与下调量
\begin{equation}
  \Delta^{+}_t=\max\!\big(G^{\mathrm{adj}}_t-G^{\mathrm{plan}}_t,\,0\big),\qquad
  \Delta^{-}_t=\max\!\big(G^{\mathrm{plan}}_t-G^{\mathrm{adj}}_t,\,0\big).
  \label{eq:q3_dpm}
\end{equation}
一天总费用为计划购电费、调整偏差费与紧急购电费之和：
\begin{equation}
  C = \sum_{t=1}^{144}\pi_t G^{\mathrm{plan}}_t
      -\dfrac{1}{2}\sum_{t=1}^{144}\pi_t \Delta^{-}_t
      +\dfrac{3}{2}\sum_{t=1}^{144}\pi_t \Delta^{+}_t
      +5\sum_{t=1}^{144}\pi_t G^{\mathrm{em}}_t.
  \label{eq:q3_cost}
\end{equation}
三项对应题面三档价格：已承诺合同按 $\pi_t$ 计价；少买“退回一半电费”（$-\tfrac12\pi_t\Delta^{-}_t$），多买按 $1.5$ 倍计价（$+\tfrac32\pi_t\Delta^{+}_t$）；真实缺口按 $5\pi_t$ 紧急购电。式~\eqref{eq:q3_cost} 是本问所有策略的统一评价指标。

\subsection{物理约束与储能动态}

以最终执行量 $G^{\mathrm{adj}}$ 替代计划量并计入紧急购电，时段电量平衡为
\begin{equation}
  G^{\mathrm{adj}}_t + P^{\mathrm{pv}}_t\Delta t + d_t + G^{\mathrm{em}}_t
  = P^{\mathrm{load}}_t\Delta t + c_t + w_t,
  \label{eq:q3_balance}
\end{equation}
储能递推与运行约束沿用前两问：
\begin{equation}
  E_t = E_{t-1}+\eta c_t - \dfrac{d_t}{\eta},\qquad
  1200\le E_t\le 10800,\qquad 0\le c_t,\,d_t\le P_{\max}\Delta t,\qquad w_t\ge 0,\qquad G^{\mathrm{em}}_t\ge 0.
  \label{eq:q3_phys}
\end{equation}
其中 $c_t,d_t$ 为交流侧充放电量，$w_t$ 为弃光，$E_t$ 为时段末储电量。这组约束\textbf{既用于优化、也用于回放}，保证“计划”与“结算”口径一致：优化时代入分位修正后的预报，回放时代入实测值。

\subsection{滚动时域线性规划与三段分位修正}

\subsubsection{$0{:}00$ 计划阶段}

每天 $0{:}00$ 用分位修正后的净负荷预报求解当天 $144$ 段线性规划，得到计划合同 $G^{\mathrm{plan}}$：
\begin{equation}
  \min_{G,c,d,w,E,G^{\mathrm{em}}}
  \sum_{j=1}^{144}\pi_j G_j + 5\sum_{j=1}^{144}\pi_j G^{\mathrm{em}}_j
  + \varepsilon\sum_{j=1}^{144}\big(c_j+d_j\big),
  \label{eq:q3_plan}
\end{equation}
约束为式~\eqref{eq:q3_balance}--\eqref{eq:q3_phys} 在修正预报下的形式。$G^{\mathrm{em}}_j$ 是 LP 内按 $5\pi$ 计价的\textbf{虚拟}紧急购电、只为保证可行性，$\varepsilon$ 是打破“同时充放”的极小正则项。

\subsubsection{$\tau>0$ 调整阶段与 48 小时展望}

在时点 $\tau\in\{6,12,18\}$，把当天已承诺合同 $g_j$ 放进“当天剩余 $+$ 次日全天”的 $48$ h 展望窗内重新决策。对当天剩余 $n_{\mathrm{today}}$ 段在 $g_j$ 上做有代价增减，对次日展望段按普通电价重新决策：
\begin{equation}
  \min_{\substack{G,\delta^{\pm},c,d,\\w,E,G^{\mathrm{em}}}}
  \sum_{j=1}^{n_{\mathrm{today}}}\Big(-\tfrac12\pi_j\delta^{-}_j+\tfrac32\pi_j\delta^{+}_j\Big)
  +\sum_{j=n_{\mathrm{today}}+1}^{n}\pi_j G_j
  +5\sum_{j=1}^{n}\pi_j G^{\mathrm{em}}_j,
  \label{eq:q3_adj}
\end{equation}
\begin{equation}
  \mathrm{s.t.}\quad G_j = g_j + \delta^{+}_j - \delta^{-}_j,\ \ \delta^{\pm}_j\ge 0\ (j\le n_{\mathrm{today}});
  \quad \text{式~\eqref{eq:q3_balance}--\eqref{eq:q3_phys}};\ \ E_0 = \mathrm{SOC}_{\tau}.
  \label{eq:q3_adj_con}
\end{equation}
每次解出后，当天剩余合同更新为新最优值，$6/12/18$ 重复即构成滚动决策（图~\ref{fig:q3_horizon}）。

\begin{figure}[!htbp]
  \centering
  \includegraphics[width=0.8\linewidth]{fig_q3_rolling.png}
  \caption{滚动时域决策示意（以 $6{:}00$ 为例）：紧邻 6 小时为锁死段（$q_L=0.60$），当天剩余为可改段（$q_O=0.50$），次日为展望段（$q_E=0.50$）}
  \label{fig:q3_horizon}
\end{figure}

\textbf{为什么拉到 48 小时？} 若只在“当天剩余”内优化，$0{:}00$ 与 $18{:}00$ 的解都会在日末把储电量耗尽、次日清晨被动；把窗延长到“当天剩余 $+$ 次日全天”后，LP 会主动为次日留储电余量，全年日末贴底从 $106$ 天压到 $1$ 天以内。次日负荷与光伏尾段用“目标日前一周同星期几”的已收盘数据，不用尚未发布的次日 $0{:}00$ 预报，避免泄漏。

\subsubsection{三段分源残差分位修正}

点预报只能给出“最可能”的净负荷，不能表达误差方向。记净负荷 $L_t=P^{\mathrm{load}}_t-P^{\mathrm{pv}}_t$，残差 $r_t=L_t-\hat L_t$ 按“同一发布时点、同一被预报钟点”分桶，得经验分布 $\hat F$。对负荷与光伏\textbf{分源}施加覆盖水平：
\begin{align}
  \tilde P^{\mathrm{load}}_t &= \max\!\Big(\hat P^{\mathrm{load}}_t + \hat F^{-1}_{\mathrm{load}}(q),\ 0\Big), \label{eq:q3_pload}\\[2pt]
  \tilde P^{\mathrm{pv}}_t  &= \max\!\Big(\hat P^{\mathrm{pv}}_t  + \hat F^{-1}_{\mathrm{pv}}(1-q),\ 0\Big). \label{eq:q3_ppv}
\end{align}
负荷取 $q$ 分位（向上抬）、光伏取 $1-q$ 分位（向下压），即“宁多备、勿少备”。分位水平按段属性取三档：锁死段 $q_L$、可改段 $q_O=0.50$、次日段 $q_E=0.50$。

\subsubsection{分位水平的报童论证}

锁死段 $q_L$ 由报童临界分位给出\textbf{搜索上界}：少买 $1$ kWh 的缺货净成本 $c_u=5\pi-\pi=4\pi$，多买 $1$ kWh 的过量成本 $c_o=\pi$，故
\begin{equation}
  q^{*}=\dfrac{c_u}{c_u+c_o}=\dfrac{4\pi}{4\pi+\pi}=0.80.
  \label{eq:q3_crit}
\end{equation}
但它只能作上界：锁死段缺额往往先由电池放电兜底、不必然触发 $5\pi$，真实 $c_u<4\pi$、临界分位应下移，故 $q_L$ 在 $\{0.50,\dots,0.75\}$ 内做一次因果回测定下（见 6.6 节）。可改段 $q_O=0.50$ 由结算对称性直接给出：多承诺净代价 $\pi-0.5\pi=0.5\pi$、少承诺净代价 $1.5\pi-\pi=0.5\pi$，两侧相等，中位数最优。

\subsection{求解与参数选择}

\textbf{（1）输入准备。} 光伏预报以发布时点对应横坐标 $0$、$T{+}k$ 对应 $k$，对 $24$ 个整点线性插值到 $10$ 分钟；插值起点取“发布前最后一个已完成时段的实测光伏”（$0{:}00$ 发布时为 $0$），避免夏季 $6/18$ 点起点硬设为零的系统低估。附件 3 的 $0{:}00$ 发布全预见期平均 MAE 约 $341$ kW、误差集中在日照时段，这正是分位余量“省不掉”的证据。

\textbf{（2）锁死段分位 $q_L$ 的一月定参。} 若在全年 $2$--$12$ 月样本上选 $q_L$，谷点在 $0.65$，但那是“看着答案挑参数”。为避免评分年选参，$q_L$ 只在 $1$ 月 $15$--$31$ 日选定：令 $\nu=\min_t \pi_t/\eta=0.4126$ 元/kWh（最低电价买入并储存 $1$ 个可用 kWh 的成本），取评分
\begin{equation}
  J(q_L)=C_{\mathrm{Jan}}(q_L)+\nu\big(6000-E^{\mathrm{Jan}}_{\mathrm{end}}(q_L)\big),
  \label{eq:q3_score}
\end{equation}
即“一月现金费用 $+$ 把月末储电相对 $6000$ kWh 的差额按 $\nu$ 折成现金”。表~\ref{tab:q3_jan} 中 $q_L=0.60$ 评分最低，故冻结 $q_L=0.60$，$2$--$12$ 月只确认、不参与挑选。

\begin{table}[!htbp]
  \centering\footnotesize
  \caption{一月（1 月 15--31 日）$q_L$ 定参网格：$q_L=0.60$ 评分最低}
  \label{tab:q3_jan}
  \begin{tabular}{@{}l r r r r@{}}
    \toprule
    $q_L$ & 一月现金费用/元 & 一月末 SOC/kWh & 评分 $J$/元 & 一月紧急/kWh \\
    \midrule
    \textbf{0.60} & 870,900.68 & 8377.37 & \textbf{869,919.88} & 3,433.65 \\
    0.63 & 871,055.16 & 8413.64 & 870,059.40 & 2,760.25 \\
    0.65 & 871,993.29 & 8440.72 & 870,986.36 & 2,310.19 \\
    0.55 & 872,521.73 & 8273.82 & 871,583.66 & 4,701.13 \\
    0.70 & 876,219.18 & 8555.60 & 875,164.86 & 1,094.51 \\
    \bottomrule
  \end{tabular}
\end{table}

\textbf{（3）两条关键杠杆的消融。} 以 $24$ h 展望 $+\beta$ 基线 $13{,}458{,}730.30$ 元为锚，全年 $334$ 天因果回测（表~\ref{tab:q3_ablation}）表明：48 h 展望单独省 $3.75$ 万元、日末贴底从 $106$ 天降到 $0$；分位修正单独省 $16.46$ 万元；二者叠加约省 $19.48$ 万元。分位修的是\textbf{晚间电量预算}缺口，不是用更多紧急电换账面费用——加入后费用与紧急电量同向变好。

\begin{table}[!htbp]
  \centering\footnotesize
  \caption{全年消融：48 h 展望与分位修正可叠加（以 24 h$+\beta$ 为锚）}
  \label{tab:q3_ablation}
  \begin{tabular}{@{}l r r r r@{}}
    \toprule
    变体 & 总费用/元 & 较锚/元 & 紧急/kWh & 贴底天数 \\
    \midrule
    24 h$+\beta$（锚） & 13,458,730.30 & 0 & 82,228.36 & 106 \\
    48 h$+\beta$ & 13,421,219.84 & $-37{,}510$ & 70,433.46 & 0 \\
    24 h$+q{=}0.65$ & 13,294,093.95 & $-164{,}636$ & 44,783.43 & 37 \\
    48 h$+q{=}0.65$ & 13,263,910.99 & $-194{,}819$ & 39,678.65 & 1 \\
    \bottomrule
  \end{tabular}

  \vspace{2pt}
  \footnotesize 注：本表仅用于分离两条杠杆的贡献，初值为 1 月滚动口径；正式数字以表~\ref{tab:q3_cmp} 的 2 月 1 日 SOC$=6000$ 口径为准。
\end{table}

\subsection{求解结果}

全年（2025-02-01 至 2025-12-31，共 $334$ 天）逐日回放，三种候选结果见表~\ref{tab:q3_cmp}。

\begin{table}[!htbp]
  \centering\footnotesize
  \caption{全年总费用：LA 最优，较不调整的 N0 省 192.2 万元}
  \label{tab:q3_cmp}
  \begin{tabular}{@{}l r r r r r@{}}
    \toprule
    策略 & 全年总费用/元 & 计划购电费/元 & 紧急购电费/元 & 紧急电量/kWh & 调整天数 \\
    \midrule
    \textbf{LA（赢家）} & \textbf{13,275,795.85} & 12,337,110.40 & 342,008.68 & 58,355.20 & 334 \\
    M0L（选择性调整） & 14,076,756.59 & 12,344,922.94 & 1,301,484.79 & 204,224.31 & 182 \\
    N0（不调整） & 15,197,639.94 & 12,342,385.85 & 2,855,254.09 & 464,543.52 & 0 \\
    \bottomrule
  \end{tabular}
\end{table}

表~\ref{tab:q3_cmp} 读出三条结论。其一，\textbf{日内调整必要}：不调整的 N0 全年 $15{,}197{,}639.94$ 元，比 LA 高 $192.2$ 万元（$+14.5\%$），紧急电量约为 LA 的 8 倍，只用 $0{:}00$ 预报扛不住 24 小时内的预报漂移。其二，\textbf{每次都调优于选择性调}：M0L 只在 $334$ 天里调整 $182$ 天，比 N0 省 $112$ 万元、仍比 LA 贵 $80$ 万元，选择性门漏掉“幅度不大但方向正确”的小修正。其三，LA 的紧急电量只有 N0 的 $12.6\%$，把风险压在计划阶段而非事后用 $5\pi$ 买单。

题目指定日期的计划购电与紧急购电见表~\ref{tab:q3_day}；储能指定日期分时段充放电量（题面表 2 格式）及完整 144 段结果见 result3.xlsx。

\begin{table}[!htbp]
  \centering\footnotesize
  \caption{指定日期计划购电量与紧急购电（题面表 1 与表 3 格式）}
  \label{tab:q3_day}
  \begin{tabular}{@{}l r r r r@{}}
    \toprule
    时间段 & 2025.3.20 & 2025.6.21 & 2025.9.23 & 2025.12.21 \\
    \midrule
    \multicolumn{5}{l}{\textbf{计划购电量 / kWh}} \\
    10:00--10:10 & 0.00 & 0.00 & 0.00 & 0.00 \\
    12:00--12:10 & 523.80 & 0.00 & 357.38 & 1006.90 \\
    14:00--14:10 & 0.00 & 0.00 & 0.00 & 0.00 \\
    16:00--16:10 & 620.65 & 55.53 & 505.11 & 836.82 \\
    18:00--18:10 & 715.88 & 411.32 & 789.37 & 719.13 \\
    20:00--20:10 & 0.00 & 0.00 & 0.00 & 0.00 \\
    全天购电/kWh & 61909.72 & 40973.80 & 66868.91 & 95888.88 \\
    全天购电费/元 & 40736.96 & 22107.38 & 45502.61 & 64039.30 \\
    \midrule
    \multicolumn{5}{l}{\textbf{紧急购电（时间段: 购电量/kWh）}} \\
    \multicolumn{5}{l}{2025.3.20：10:00--10:10 31.93；20:40--20:50 131.05} \\
    \multicolumn{5}{l}{2025.6.21：无紧急购电} \\
    \multicolumn{5}{l}{2025.9.23：8:50--10:00 446.61；20:30--20:50 143.22；21:00--21:10 60.15} \\
    \multicolumn{5}{l}{2025.12.21：8:00--8:20 23.68；10:40--11:00 254.38；20:40--20:50 43.38} \\
    \bottomrule
  \end{tabular}
\end{table}

夏季代表日（2025-06-21）滚动执行过程见图~\ref{fig:q3_dispatch}：净负荷晚高峰 $4004.9$ kW@$19.2$ h 由储能放电削峰，日内调整较计划少买 $127$ kWh。

\begin{figure}[!htbp]
  \centering
  \includegraphics[width=0.7\linewidth]{fig_q3_dispatch_0621.png}
  \caption{夏季代表日（2025-06-21）滚动购电策略：储能晚高峰放电削峰，调整购电较计划少买 127 kWh}
  \label{fig:q3_dispatch}
\end{figure}

\subsection{模型检验与稳健性}

\textbf{（1）独立复算。} 导出 result3.xlsx 后由独立脚本 \texttt{check\_result3.py} 从单元格与附件 1 电价重算逐日合同费、偏差结算、SOC 平衡与跨日连续，不复用滚动代码：正式结果恰好 $334$ 天，储电量全程 $[1200,10800]$、功率不超 $5000$ kW，$2$ 月 $1$ 日 $0{:}00$ 为 $6000$ kWh、相邻日首尾相接，全程无泄漏错误。

\textbf{（2）与先知下界的距离。} 阶段 1（2 月 1--14 日）先知（用实测净负荷解单日 LP 的完美信息下界）总费用 $585{,}705.84$ 元，LA 为 $636{,}942.33$ 元，仅高 $8.7\%$；同期 N0 为 $653{,}094.43$ 元。LA 已逼近先知下界，剩余差距主要来自预报残差而非模型结构。

\textbf{（3）参数稳健性。} $q_L$ 在 V 形谷底附近不敏感：$0.60$ 与 $0.65$ 全年只差约 $0.08\%$，不存在“刀尖上的最优”。一月定参 $q_L=0.60$ 在 $2$--$12$ 月确认重跑中通过门禁、与因果重跑逐位一致。

\textbf{（4）失效模式诊断。} 对旧 $24$ h$+\beta$ 方案逐段回放表明，紧急购电 $88.8\%$ 来自\textbf{晚间储能电量耗尽}而非功率触顶，$20{:}00$ 单小时占 $47.4\%$。这解释了“加大光伏余量”与“回放留 SOC”为何无效——都没碰到“晚间电量预算”这个真约束；新方案把修正做进 LP 本身，日末贴底降到 $1$ 天。

\subsection{本问结论}

本问需要引入 $6{:}00$、$12{:}00$、$18{:}00$ 时点预报调整购电：不调整的 N0 全年 $15{,}197{,}639.94$ 元，比最优策略 LA 高 $192.2$ 万元（$+14.5\%$），说明日内调整必要；每次都调优于选择性调（M0L 比 LA 贵 $80$ 万元）。最优策略为 LA——$48$ h 展望加三段分源残差分位修正（$q_L=0.60$，一月定参），全年总费用 $13{,}275{,}795.85$ 元，紧急电量较 N0 下降 $87.4\%$。完整策略见 result3.xlsx。
